\section{Proof sketches}

This section provides proof sketches for the falsifiability criteria stated
in \Cref{sec:falsifiability-criteria}.
All arguments are purely structural and operate exclusively on the
existence or non-existence of compatible states in the formal state space
\(\Omega(K_{EA})\).
No assumptions about empirical interpretation, diagnostic heuristics, or
procedural completeness are involved.

\medskip

\begin{proof}[Proof sketch for \Cref{thm:necessary-falsifiability}]
Assume that there exists at least one model realization
\(s \in \Omega(K_{EA})\) such that the admissible observation \(O\) is
compatible with \(s\), i.e., all statements in \(O\) can be embedded into
\(s\) without violating any mandatory invariant, threshold, or structural
constraint.
By definition, internal falsification requires the absence of any
compatible model realization.
The existence of such an \(s\) therefore excludes falsification.
\end{proof}

\medskip

\begin{proof}[Proof sketch for \Cref{thm:sufficient-falsifiability}]
Assume that an admissible observation \(O\) logically implies inconsistency
with at least one invariant, threshold, or structural constraint that is
mandatory for all admissible model realizations
\(s \in \Omega(K_{EA})\).
By definition, every admissible realization must satisfy all mandatory
constraints.
If \(O\) implies inconsistency with any such constraint, then no admissible
realization can exist that is compatible with \(O\).
The compatibility set is therefore empty, and internal falsification of the
enterprise-as-continuum model follows.
\end{proof}

\medskip

\begin{proof}[Proof sketch for \Cref{prop:nonfalsifying}]
Non-identifiability or underdetermination means that multiple distinct model
realizations
\[
s_1, s_2, \dots \in \Omega(K_{EA})
\]
remain compatible with the same admissible observation \(O\).
This situation reflects a multiplicity of admissible realizations rather
than a logical contradiction.
Since at least one compatible realization exists, the set of compatible
states is non-empty and falsification does not occur.
\end{proof}

\medskip

\begin{proof}[Proof sketch for \Cref{prop:stop-not-falsification}]
A \emph{STOP} outcome asserts that no further admissible refinement or
continuation of the diagnostic process is possible without violating model
constraints.
It is therefore a statement about the limits of admissible inference under
the given observation interface.
Crucially, STOP does not assert that the state space
\(\Omega(K_{EA})\) is empty, nor that no compatible model realization
exists.
Hence a STOP outcome cannot imply falsification.
\end{proof}

\medskip

\begin{remark}
All proofs reduce to a single meta-criterion.
Let
\[
\Omega_O \;=\; \{\, s \in \Omega(K_{EA}) \mid s \text{ is compatible with }
O \,\}
\]
denote the compatibility set induced by an admissible observation \(O\).
Internal falsification is equivalent to the condition
\(\Omega_O = \varnothing\).
The enterprise-as-continuum model is therefore falsifiable precisely to the
extent that admissible observations can, in principle, force the emptiness
of \(\Omega_O\) by inducing internal inconsistency with mandatory
constraints.
All other outcomes correspond to non-empty compatibility sets and are
explicitly non-falsifying.
\end{remark}
